% ---
% RESUMOS
% ---

% RESUMO em português
\renewcommand{\resumoname}{Objetivos e Resumo}
\setlength{\absparsep}{18pt} % ajusta o espaçamento dos parágrafos do resumo
\begin{resumo}

\lipsum[3]

% \hspace*{1.3cm}No primeiro capítulo é apresentada uma breve introdução, estabelecem-se as principais classificações de digitalizadores 3D, subdividindo-os em seguida. Por fim, são apresentadas as principais aplicações dos aparelhos de escaneamento tridimensional.

% \hspace*{1.3cm}No segundo capítulo é apresentado o estado da arte dos digitalizadores, contendo artigos de revista e conferência, servindo de guia para o acontece de inovação no mundo nessa área.

% \hspace*{1.3cm}No terceiro capítulo tem-se que alguns aparelhos comerciais são apresentados e uma breve comparação é feita entre eles.

% \hspace*{1.3cm}\textcolor{red}{Esse documento ainda está em formação, portanto, mais elementos serão acrescentados ao longo do projeto. Espera-se que ao longo do tempo de projeto tenha-se a possibilidade de constituir publicações e mais relatórios, com a ajuda deste trabalho.}

 \textbf{Palavras-chaves}: \textit{template}, \textit{report}, \textit{relatório}.
\end{resumo}

% ABSTRACT in english
%\begin{resumo}[Abstract]
% \begin{otherlanguage*}{english}
%   This is the english abstract.
%
%   \vspace{\onelineskip}
% 
%   \noindent 
%   \textbf{Keywords}: latex. abntex. text editoration.
% \end{otherlanguage*}
%\end{resumo}